Parallel torsion gives a setting in which an integral balanced trace becomes a pointwise constraint. Let $g$ be a Hermitian metric on the connected compact manifold $X$, of dimension $n\geq2$, and let $D$ be a real connection on $TX$ satisfying
\[
 Dg=0,\qquad DJ=0,\qquad DT^D=0,
 \qquad T^D(U,V)=D_UV-D_VU-[U,V].
\]
The balanced and pluriclosed metrics considered below need not equal $g$. Write
\[
 E_D^k=\{\alpha\in\Omega^k(X,\mathbb R):D\alpha=0\}.
\]
Evaluation at one point is injective on this space, so it is finite-dimensional. Its complexification inherits the bidegree decomposition because $DJ=0$.

\begin{lemma}\label{lem:parallel-positive-projection}\index{Parallel torsion!positive projection}
There is a linear projection $S:\Omega^\bullet(X,\mathbb R)\to E_D^\bullet$ with $S\mathrm d=\mathrm dS$. Its complexification preserves bidegree and commutes with $\partial$ and $\bar\partial$. It preserves strict positivity in bidegrees $(1,1)$ and $(n-1,n-1)$.
\end{lemma}
\begin{proof}
For a parallel form, the formula for exterior differentiation with a connection expresses $\mathrm d\alpha$ as an algebraic contraction of $T^D$ with $\alpha$. Thus $DT^D=0$ implies $\mathrm dE_D^k\subseteq E_D^{k+1}$. The Hodge star of $g$ is also parallel and maps $E_D^k$ onto $E_D^{2n-k}$. Consequently the pairing
\[
 E_D^k\times E_D^{2n-k}\longrightarrow\mathbb R,
 \qquad (\alpha,\beta)\longmapsto\int_X\alpha\wedge\beta
\]
is nondegenerate: pairing a nonzero $\alpha$ with $*_g\alpha$ gives its positive squared $L^2$ norm. There is a unique $S_k\alpha\in E_D^k$ such that
\begin{equation}\label{eq:parallel-projection-pairing}
 \int_X S_k\alpha\wedge\beta=\int_X\alpha\wedge\beta
 \qquad(\beta\in E_D^{2n-k}).
\end{equation}
This defines a projection. For $\gamma\in E_D^{2n-k-1}$, Stokes and $\mathrm d\gamma\in E_D^{2n-k}$ give
\[
 \int_X\mathrm dS_k\alpha\wedge\gamma
 =(-1)^{k+1}\int_X\alpha\wedge\mathrm d\gamma
 =\int_X S_{k+1}(\mathrm d\alpha)\wedge\gamma.
\]
Both forms being compared are parallel. Nondegeneracy proves $\mathrm dS_k\alpha=S_{k+1}\mathrm d\alpha$.

To check positivity, fix $p\in X$. Let $H$ be the closure of the full $D$-holonomy group at $p$, a compact subgroup of the unitary group, with probability Haar measure $m_H$. For a path from $p$ to $x$, write $P_x:T_pX\to T_xX$ for parallel transport. The average
\[
 R_\alpha(x)=\int_H h^*(P_x^*\alpha_x)\,\mathrm dm_H(h)
\]
is independent of the path, since a change of path acts by an element of $H$. It is a smooth $H$-invariant tensor at $p$: paths can be chosen smoothly on each coordinate neighborhood. Average it over $X$,
\[
 a_p=\frac1{\operatorname{Vol}_g(X)}
                \int_X R_\alpha(x)\,\mathrm dV_g(x).
\]
An $H$-invariant tensor extends uniquely to a global parallel form $\widetilde S_k\alpha$. For parallel $\beta$, write $\alpha\wedge\beta=f\,\mathrm dV_g$. Unitary transport and holonomy averaging give
\[
 R_\alpha(x)\wedge\beta_p=f(x)\,\mathrm dV_{g,p}.
\]
Integrating this equality shows that $\widetilde S_k\alpha$ satisfies \eqref{eq:parallel-projection-pairing}; hence $\widetilde S_k=S_k$.

Every transport and holonomy action preserves reality, bidegree and the positive cones. A smooth strictly positive $(q,q)$-form, for $q=1$ or $n-1$, is bounded below by a positive multiple of the corresponding power of the form of $g$, uniformly on the compact manifold. Both probability averages preserve that lower bound. Therefore $S$ preserves strict positivity and bidegree. Comparing bidegrees in $S\mathrm d=\mathrm dS$ proves commutation with $\partial$ and $\bar\partial$.
\end{proof}

The final step of the construction uses the same-metric implication mentioned above.
\begin{lemma}\label{lem:balanced-pluriclosed-closed}
If a Hermitian form $\zeta$ satisfies $\mathrm d\zeta^{n-1}=0$ and $\ddc\zeta=0$, then $\mathrm d\zeta=0$.
\end{lemma}
\begin{proof}
For $n=2$, the balanced condition already states $\mathrm d\zeta=0$. Suppose $n\geq3$ and set $\alpha=\partial\zeta$. Balancedness gives $\alpha\wedge\zeta^{n-2}=0$, so $\alpha$ is a primitive $(2,1)$-form. Differentiating the balanced power in both complex directions gives
\[
 0=\partial\bar\partial\zeta^{n-1}
  =(n-1)\partial\bar\partial\zeta\wedge\zeta^{n-2}
    +(n-1)(n-2)\alpha\wedge\bar\alpha\wedge\zeta^{n-3}.
\]
Pluriclosedness cancels the first term. Pointwise Hermitian linear algebra for a primitive $(2,1)$-form gives
\[
 -\mathrm i\,\alpha\wedge\bar\alpha\wedge
                 \frac{\zeta^{n-3}}{(n-3)!}
       =|\alpha|_\zeta^2\frac{\zeta^n}{n!}.
\]
Thus $\alpha=0$, and reality of $\zeta$ gives $\mathrm d\zeta=0$.
\end{proof}

\begin{theorem}\label{thm:parallel-torsion-kahler}\index{Parallel torsion!Kahler metric@K\"ahler metric}\index{Balanced metric!parallel torsion}\index{Pluriclosed metric!parallel torsion}
Under the stated hypotheses on $D$, if $X$ admits a balanced metric and a pluriclosed metric, then it admits a $D$-parallel K\"ahler metric.
\end{theorem}
\begin{proof}
Let $\beta$ be balanced and $\omega$ pluriclosed. By \cref{lem:parallel-positive-projection}, $S(\beta^{n-1})$ is strictly positive, parallel and closed. It has a unique smooth positive $(n-1)$-st root $B$. Indeed, after the fixed wedge normalization, the matrix map is $G\mapsto C=(\det G)G^{-1}$, whose inverse is
\[
 G=(\det C)^{1/(n-1)}C^{-1}.
\]
Uniqueness makes the root invariant under parallel transport. Thus $DB=0$ and $\mathrm dB^{n-1}=0$. Similarly $P=S\omega$ is a parallel positive pluriclosed form.

Consider the finite-dimensional space of real parallel exact-type directions
\[
 \mathcal A_D=\{(\mathrm d\xi)^{1,1}:\xi\in E_D^1\}.
\]
Every such direction is $\partial\bar\partial$-closed, since
\[
 (\mathrm d\xi)^{1,1}=\partial\xi^{0,1}+\bar\partial\xi^{1,0}.
\]
For a parallel positive form $Q$, the volume functional has derivative
\begin{equation}\label{eq:parallel-volume-variation}
 \left.\frac{\mathrm d}{\mathrm dt}\right|_{t=0}
       \frac1{n!}\int_X(Q+t(\mathrm d\xi)^{1,1})^n
 =\frac1{(n-1)!}\int_X\xi\wedge\mathrm dQ^{n-1}.
\end{equation}
The $(2,0)$ and $(0,2)$ parts of $\mathrm d\xi$ pair to zero with $Q^{n-1}$. Since $\mathrm dQ^{n-1}$ is parallel, nondegeneracy of the parallel pairing shows that all these derivatives vanish exactly when $Q$ is balanced.

Apply this to $B$. For $a\in\mathcal A_D$, the function $\operatorname{tr}_B a$ is parallel and hence constant. Its integral against $B^n$ vanishes by \eqref{eq:parallel-volume-variation}, so it vanishes pointwise. Therefore the positive affine slice
\[
 \mathcal P=\{Q\in P+\mathcal A_D:Q>0\}
\]
has the fixed pointwise trace $\operatorname{tr}_B Q=\operatorname{tr}_B P=:C>0$.

Evaluate the parallel forms at one point. The positive eigenvalues of $Q$ relative to $B$ sum to $C$, so they are bounded above by $C$. Thus $\mathcal P$ is bounded in its finite-dimensional affine space. Its closure is precisely
\[
 \overline{\mathcal P}=\{Q\in P+\mathcal A_D:Q\geq0\},
\]
because $(1-t)Q+tP$ is positive for $t>0$ whenever $Q$ is semipositive. This closure is compact. The relative eigenvalues and determinant are constant on $X$, since both forms are parallel. Hence
\[
 \frac1{n!}\int_XQ^n
       =\det(B^{-1}Q)\frac1{n!}\int_XB^n.
\]
At a relative boundary point the semipositive matrix is singular, so this volume is zero. It is positive at $P$. A volume maximum on the compact closure therefore lies in $\mathcal P$.

At that interior maximum all derivatives in \eqref{eq:parallel-volume-variation} vanish. The maximizing form is balanced, and it is pluriclosed because it lies in $P+\mathcal A_D$. By \cref{lem:balanced-pluriclosed-closed} it is closed. It is already positive and $D$-parallel, which proves the theorem.
\end{proof}

The trace constraint in this argument is pointwise because the metrics and the variation directions are parallel. This is the additional structure that makes the positive slice compact and its boundary volume vanish; the conserved balanced trace for a general family does not supply it.
